\documentclass[12pt]{article}

\usepackage[brazilian]{babel}
\usepackage[utf8]{inputenc}
\usepackage[T1]{fontenc}
\usepackage{graphicx}
\usepackage{geometry}
\usepackage{booktabs}
\usepackage{float}
\usepackage{longtable}
\usepackage[table]{xcolor}
\usepackage{tcolorbox}
\usepackage{fancyvrb}
\usepackage{hyperref}

\geometry{a4paper, left=3cm, top=3cm, right=2cm, bottom=2cm}
\definecolor{rosa}{HTML}{D6336C}
\hypersetup{colorlinks=true, linkcolor=rosa, urlcolor=rosa}
\setlength{\parindent}{0pt}
\setlength{\parskip}{6pt}
\newtcolorbox{dica}[1][Atenção]{colback=rosa!6, colframe=rosa, title=#1, fonttitle=\bfseries, boxrule=0.6pt}

\begin{document}
\begin{titlepage}
    \centering
    \textsc{\Large Universidade Cruzeiro do Sul}\\[0.5cm]
    \textsc{\Large Engenharia de Software}\\[0.5cm]
    \textsc{Projeto Integrador Transdisciplinar II}
    \vspace{2.5cm}

    \hrule \vspace{0.4cm}
    {\huge \bfseries Documento de Testes\\[0.3cm] Loja Virtual de Cupcakes}\\[0.4cm]
    \hrule
    \vspace{3cm}

    \begin{flushleft}
        \Large \emph{Aluno:}\\
        Fellipe Augusto Gonçalves Magalhães
    \end{flushleft}
    \vfill
    {\large Novembro de 2026}
\end{titlepage}

\tableofcontents
\newpage

% ---------------------------------------------------------------- 1
\section{Introdução}

Este documento descreve os testes automatizados do sistema \textbf{Loja Virtual de Cupcakes}, os resultados obtidos, os defeitos encontrados e corrigidos e a rastreabilidade entre os testes e as histórias de usuário (US01 a US15) e regras de negócio (RN01 a RN15) definidas no Projeto Integrador I.

O objetivo é demonstrar, de forma repetível, que o comportamento esperado foi realizado. Os testes cobrem a \textbf{verificação} (o sistema foi construído conforme a especificação) e preparam a \textbf{validação} com usuários, etapa seguinte do projeto (seção~\ref{sec:limites}).

\begin{table}[H]
\centering\renewcommand{\arraystretch}{1.3}
\begin{tabular}{|p{6cm}|p{7.4cm}|}
\hline
\rowcolor{gray!30} \textbf{Indicador} & \textbf{Resultado} \\ \hline
Testes automatizados de back-end e banco (JUnit 5) & 131 casos, todos aprovados \\ \hline
Testes automatizados de interface (Playwright) & 31 casos, todos aprovados \\ \hline
Histórias de usuário com ao menos um teste & 15 de 15 \\ \hline
Regras de negócio com ao menos um teste & 15 de 15 \\ \hline
Defeitos encontrados e corrigidos & 2 (seção~\ref{sec:defeitos}) \\ \hline
\end{tabular}
\caption{Resumo geral}
\end{table}

% ---------------------------------------------------------------- 2
\section{Estratégia de testes}

Os testes seguem uma pirâmide: muitos testes rápidos nas camadas internas e menos testes, mais amplos, nas camadas externas.

\begin{table}[H]
\centering\small\renewcommand{\arraystretch}{1.3}
\begin{tabular}{|p{3.3cm}|p{5.3cm}|p{5.3cm}|}
\hline
\rowcolor{gray!30} \textbf{Nível} & \textbf{O que verifica} & \textbf{Como} \\ \hline
Unitário & Funções isoladas: conversão JSON, hash de senha, sessões & JUnit 5, sem banco \\ \hline
Integração (DAO + banco) & Cada consulta e transação do Model contra o PostgreSQL real & JUnit 5; o banco é recriado antes de cada teste \\ \hline
Banco de dados & Restrições, índices únicos, triggers, views e permissões & JUnit 5 com SQL direto \\ \hline
Ponta a ponta (API) & Rotas HTTP, perfis, validações e fluxos completos (Controller + DAO + banco) & JUnit 5 sobe o servidor em uma porta livre e faz requisições HTTP \\ \hline
Interface (View) & Telas, mensagens e navegação no navegador & Playwright (Chromium), simulando o usuário \\ \hline
\end{tabular}
\caption{Níveis de teste}
\end{table}

\textbf{Critérios adotados em cada caso:} um caso de teste verifica um comportamento esperado e tem resultado objetivo (aprovado ou reprovado). Cada caso traz a indicação da história (USxx) e da regra de negócio (RNxx) a que se refere. Os casos incluem o caminho feliz, as entradas inválidas e os limites (por exemplo, quantidade 0, 1 e 100).

\textbf{Isolamento:} antes de cada teste de back-end, o esquema do banco de testes é apagado e recriado pelo mesmo script SQL do projeto, com os dados iniciais. Assim, nenhum teste depende do resultado de outro, e o próprio script de criação do banco é exercitado a cada execução.

% ---------------------------------------------------------------- 3
\section{Ambiente e como executar}

\subsection{Ambiente usado}

\begin{itemize}
\item Java 21 (exige-se JDK 17 ou superior), JUnit 5 em modo console (\texttt{junit-platform-console-standalone}), driver JDBC PostgreSQL 42.7.2.
\item PostgreSQL 16, em um banco separado do banco de uso (\texttt{loja\_cupcakes\_test}).
\item Interface: Python 3, Playwright e Chromium.
\end{itemize}

\begin{dica}[Proteção do banco]
Os testes apagam e recriam o esquema \texttt{loja}. Por segurança, só rodam se o nome do banco contiver ``test''; caso contrário, o teste aborta antes de tocar em qualquer dado.
\end{dica}

\subsection{Executar os testes de back-end}

\begin{enumerate}
\item Crie o banco de testes (com o container do banco já ligado):
\begin{Verbatim}[fontsize=\footnotesize, frame=single]
docker exec -it loja-db psql -U postgres -c "CREATE DATABASE loja_cupcakes_test"
\end{Verbatim}
\item Na pasta do projeto, rode (\texttt{testar.sh} no Linux/Mac):
\begin{Verbatim}[fontsize=\footnotesize, frame=single]
testar.bat
\end{Verbatim}
\end{enumerate}

O script compila o sistema e os testes e executa tudo. O relatório em XML fica na pasta \texttt{relatorios}. Se o banco usar outra senha ou porta, defina \texttt{TEST\_DB\_URL}, \texttt{TEST\_DB\_USER} e \texttt{TEST\_DB\_PASSWORD}.

\subsection{Executar os testes de interface}

\begin{enumerate}
\item Instale uma vez: \texttt{pip install playwright} e \texttt{playwright install chromium}.
\item Recrie o banco de testes com os dados iniciais e suba o sistema apontando para ele (\texttt{DB\_URL} com o banco \texttt{loja\_cupcakes\_test}).
\item Rode: \texttt{python testes/interface/testar\_interface.py http://localhost:8080}.
\end{enumerate}

O script grava \texttt{relatorios/interface.json} e salva uma imagem da tela de cada teste reprovado.

% ---------------------------------------------------------------- 4
\section{Resultados}

\subsection{Back-end e banco de dados}

\begin{table}[H]
\centering\small\renewcommand{\arraystretch}{1.25}
\begin{tabular}{|p{4.6cm}|p{3.2cm}|p{3.7cm}|c|c|}
\hline
\rowcolor{gray!30} \textbf{Classe de teste} & \textbf{Nível} & \textbf{Foco} & \textbf{Casos} & \textbf{Aprovados} \\ \hline
\texttt{JsonTest} & Unitário & Conversor JSON & 10 & 10 \\ \hline
\texttt{SenhaESessaoTest} & Unitário & Hash de senha e sessões & 5 & 5 \\ \hline
\texttt{UsuarioEEnderecoDaoTest} & Integração (DAO + banco) & Usuários, login e endereços & 13 & 13 \\ \hline
\texttt{CupcakeDaoTest} & Integração (DAO + banco) & Vitrine, filtro e produtos & 15 & 15 \\ \hline
\texttt{CarrinhoDaoTest} & Integração (DAO + banco) & Carrinho de compras & 11 & 11 \\ \hline
\texttt{PedidoDaoTest} & Integração (DAO + banco) & Pedido, pagamento, status e dashboard & 29 & 29 \\ \hline
\texttt{BancoDeDadosTest} & Banco de dados & Restrições, triggers, views e permissões & 13 & 13 \\ \hline
\texttt{ApiTest} & Ponta a ponta (API HTTP) & Rotas, perfis, validações e fluxos completos & 35 & 35 \\ \hline
\textbf{Total} & & & \textbf{131} & \textbf{131} \\ \hline
\end{tabular}\caption{Resumo dos testes automatizados de back-end e banco (JUnit 5)}\end{table}

Saída resumida do JUnit na última execução (todos os casos aprovados, tempo total de cerca de 21 segundos):

\begin{Verbatim}[fontsize=\footnotesize, frame=single]
Test run finished after 20596 ms
[ 11 containers found      ]
[ 0 containers skipped    ]
[ 11 containers started    ]
[ 0 containers aborted    ]
[ 11 containers successful ]
[ 0 containers failed     ]
[ 138 tests found           ]
[ 0 tests skipped         ]
[ 138 tests started         ]
[ 0 tests aborted         ]
[ 138 tests successful      ]
[ 0 tests failed          ]
\end{Verbatim}

\subsection{Interface}

Os 31 casos de interface foram executados em um navegador Chromium automatizado, contra o sistema real e o banco de testes, e todos foram aprovados. Eles cobrem: tela inicial só com login, cadastro e login, vitrine e filtro, detalhes, carrinho na aba lateral, finalização do pedido, pagamento recusado e aprovado, Meus pedidos, painel do dono (dashboard, status, cadastro e edição de produto) e proteção contra HTML malicioso no nome de produto. A lista completa está no Apêndice B.

Para garantir que a verificação detecta erros, foi feito um controle negativo: uma checagem propositalmente falsa (esperar a loja visível antes do login) foi reprovada pelo mesmo mecanismo.

% ---------------------------------------------------------------- 5
\section{Defeitos encontrados e corrigidos}
\label{sec:defeitos}

\begin{table}[H]
\centering\small\renewcommand{\arraystretch}{1.3}
\begin{tabular}{|c|p{4.0cm}|p{3.5cm}|p{3.5cm}|}
\hline
\rowcolor{gray!30} \textbf{ID} & \textbf{Defeito} & \textbf{Correção} & \textbf{Teste de regressão} \\ \hline
D1 & Enviar um JSON cortado (por exemplo, \texttt{\{"email":}) fazia o servidor responder erro 500 (``Erro interno'') em vez de erro 400, porque o leitor lançava \texttt{String\-Index\-Out\-Of\-Bounds\-Exception} & \texttt{Json.\-lerObjeto} passou a converter o erro de leitura em \texttt{Illegal\-Argument\-Exception}, que a API responde com 400 & \texttt{Json\-Test}: JSON malformado; \texttt{Api\-Test}: JSON malformado devolve 400 \\ \hline
D2 & Enviar número decimal em campo inteiro (por exemplo, quantidade 1,5) gerava \texttt{Arithmetic\-Exception} e erro 500 & \texttt{Json.longo} passou a lançar \texttt{Illegal\-Argument\-Exception} com a mensagem ``O campo ... deve ser um número inteiro'' & \texttt{Json\-Test}: decimal em inteiro; \texttt{Api\-Test}: quantidade inválida (1,5) \\ \hline
D3 & A tela de login continuava visível depois do acesso, porque o estilo \texttt{\#tela-login \{ display:flex \}} anulava a classe que a oculta & Regra \texttt{.oculto \{ display:none !important \}} em \texttt{style.css} & Caso de interface UI01 (tela inicial só com o login) \\ \hline
D4 & (Colegas G. Siqueira e E. Peron) Produto com estoque zero e quantidade maior que o estoque podiam ser colocados no carrinho; o erro só aparecia no fechamento do pedido & \texttt{Carrinho\-Dao} recusa produto esgotado e quantidade (somada ao que já está no carrinho) acima do estoque, com mensagem clara; a vitrine marca ``Esgotado'' e desabilita o botão & \texttt{Carrinho\-Dao\-Test}: esgotado, acima do estoque; \texttt{Api\-Test}: 400 com mensagem, estoque na vitrine; UI32, UI35 e UI36 \\ \hline
D5 & (Colega E. Peron) Não era possível escolher a quantidade de uma vez: era preciso adicionar uma unidade por clique & Campo de quantidade na vitrine e nos detalhes; no carrinho, botões + e − e campo numérico (nova rota \texttt{PUT /api/carrinho/itens/\{id\}}) & \texttt{Carrinho\-Dao\-Test}: definir quantidade; \texttt{Api\-Test}: PUT; UI33 e UI34 \\ \hline
D6 & (Colega W. Spagnol) O dashboard do dono não atualizava sozinho & Atualização automática a cada 10 s (pausa com a aba oculta), horário da última atualização e botão ``Atualizar agora'' & UI37 e UI38 \\ \hline
\end{tabular}
\caption{Defeitos encontrados nos testes}
\end{table}

O D1 foi detectado pela execução dos testes (quatro casos reprovaram na primeira execução; dois deles eram erros na escrita dos próprios testes, que usavam um leitor de JSON que não lê listas, e foram ajustados). O D2 foi identificado na análise do mesmo trecho de código e recebeu testes de regressão antes da última execução completa. Após as correções, todos os casos foram aprovados. Os defeitos D4, D5 e D6 foram apontados pelos colegas na rodada de testes de 28/10 a 02/11/2026 (relato completo no laudo de qualidade entregue com o projeto); para cada um foi escrito primeiro o teste de regressão, que reprovou com o código antigo (por exemplo, a API respondia 200 ao adicionar um produto esgotado), e depois a correção.

% ---------------------------------------------------------------- 6
\section{Rastreabilidade}

\subsection{Histórias de usuário}

A tabela mostra quantos casos de teste referenciam cada história. Um caso pode referenciar mais de uma.

\begin{longtable}{|p{1.2cm}|p{4.6cm}|c|c|c|p{2.8cm}|}
\hline
\rowcolor{gray!30} \textbf{ID} & \textbf{História} & \textbf{Back-end} & \textbf{Interface} & \textbf{Total} & \textbf{Resultado} \\ \hline
\endhead
US01 & Visualizar vitrine & 5 & 2 & 7 & Aprovados \\ \hline
US02 & Filtrar por categoria & 6 & 1 & 7 & Aprovados \\ \hline
US03 & Ver detalhes & 3 & 1 & 4 & Aprovados \\ \hline
US04 & Adicionar ao carrinho & 9 & 1 & 10 & Aprovados \\ \hline
US05 & Remover do carrinho & 3 & 2 & 5 & Aprovados \\ \hline
US06 & Visualizar carrinho & 4 & 1 & 5 & Aprovados \\ \hline
US07 & Endereço de entrega & 8 & 3 & 11 & Aprovados \\ \hline
US08 & Forma de pagamento & 2 & 0 & 2 & Aprovados \\ \hline
US09 & Confirmar pedido & 13 & 2 & 15 & Aprovados \\ \hline
US10 & Processar pagamento & 8 & 1 & 9 & Aprovados \\ \hline
US11 & Acompanhar pedido & 4 & 1 & 5 & Aprovados \\ \hline
US12 & Cadastrar produto & 6 & 2 & 8 & Aprovados \\ \hline
US13 & Editar produto & 5 & 1 & 6 & Aprovados \\ \hline
US14 & Atualizar status & 8 & 2 & 10 & Aprovados \\ \hline
US15 & Dashboard de vendas & 6 & 1 & 7 & Aprovados \\ \hline
\end{longtable}

\begin{dica}[Histórias parcialmente implementadas]
Alguns pontos do planejamento não estão implementados (por exemplo, consulta de CEP no ViaCEP, \textit{webhook} do gateway e gráficos do dashboard). Os testes verificam o comportamento que existe hoje; esses pontos estão listados no capítulo de implementação do documento do projeto e entram nos testes quando forem implementados.
\end{dica}

\subsection{Regras de negócio}

\begin{longtable}{|c|c|p{9.4cm}|}
\hline
\rowcolor{gray!30} \textbf{Regra} & \textbf{Casos} & \textbf{Verificada por (exemplos)} \\ \hline
\endhead
RN01 & 2 & \small Produto inativo não aparece na vitrine da API; Vitrine lista somente produtos ativos \\ \hline
RN02 & 2 & \small Categoria inválida gera erro 400; Categoria inexistente retorna lista vazia \\ \hline
RN03 & 2 & \small Dono edita preço e estoque (detalhes mostram o dado novo); desativar tira da vitrine; Editar produto altera preço, estoque e categorias (dado atualizado) \\ \hline
RN04 & 4 & \small Quantidade fora de 1 a 99 é recusada com 400; Definir a quantidade de um item substitui o valor (D5) \\ \hline
RN05 & 1 & \small Quantidade negativa é rejeitada \\ \hline
RN06 & 1 & \small Carrinho calcula subtotal por item \\ \hline
RN07 & 4 & \small Endereço de outro cliente ou inexistente é recusado; CEP inválido e UF em minúsculas são rejeitados pelo banco \\ \hline
RN08 & 2 & \small Forma de pagamento inválida é rejeitada e nada é gravado; Forma de pagamento inválida devolve 400 \\ \hline
RN09 & 11 & \small Produto esgotado não pode ser adicionado ao carrinho (D4); Quantidade acima do estoque é rejeitada, inclusive somada ao que já está no carrinho (D4) \\ \hline
RN10 & 3 & \small Pagamento aprovado muda o pedido para PAGO; Pagamento processado precisa ter data de processamento \\ \hline
RN11 & 2 & \small Status inexistente e pedido inexistente geram erro; Status inexistente, ausente ou pedido inexistente devolvem 400 \\ \hline
RN12 & 3 & \small Nome de produto repetido é rejeitado; Cadastro recusa campos obrigatórios vazios, preço zero e nome repetido \\ \hline
RN13 & 2 & \small Editar produto inexistente gera erro; Editar produto inexistente devolve 400 \\ \hline
RN14 & 6 & \small Pedido entregue não pode mais mudar de status; Dono muda o status seguindo o fluxo e é barrado ao pular etapas \\ \hline
RN15 & 1 & \small Vendas só contam pedidos já pagos e refletem o dado atual \\ \hline
\end{longtable}

% ---------------------------------------------------------------- 7
\section{Requisitos não funcionais}

\begin{table}[H]
\centering\small\renewcommand{\arraystretch}{1.3}
\begin{tabular}{|c|p{4.6cm}|p{6.2cm}|}
\hline
\rowcolor{gray!30} \textbf{ID} & \textbf{Requisito} & \textbf{Situação nos testes} \\ \hline
RNF01, RNF02 & Vitrine em menos de 2 s e filtro em menos de 1 s & Verificado: ambos respondem em menos de 1 s (teste de desempenho da API) \\ \hline
RNF03 & Interface responsiva & Não verificado automaticamente (a interface usa grade flexível; falta teste em celular) \\ \hline
RNF04 & Persistência do carrinho & Verificado: o carrinho fica no banco (testes do carrinho e recarregamento da página) \\ \hline
RNF08, RNF12 & Segurança & Verificado: perfis (401 e 403), hash de senha, permissões do banco, escape de HTML \\ \hline
RNF09, RNF13 & Consistência e integridade & Verificado: transação do pedido, restrições e chaves do banco \\ \hline
RNF05, RNF06, RNF11, RNF14 & Atualização instantânea e tempo real & Parcial: o carrinho atualiza na hora; o status dos pedidos usa consulta periódica \\ \hline
RNF10, RNF15 & Disponibilidade e desempenho do dashboard & Não verificado (exigiria teste de carga) \\ \hline
\end{tabular}
\caption{Requisitos não funcionais e verificação}
\end{table}

% ---------------------------------------------------------------- 8
\section{Limites desta etapa e próximos passos}
\label{sec:limites}

\begin{itemize}
\item \textbf{Validação com usuários:} os testes automatizados verificam a especificação. A rodada com cinco colegas foi feita entre 28/10 e 02/11/2026 e gerou os defeitos D4, D5 e D6, já corrigidos. Os colegas também apontaram que não existe recuperação de senha e que o pagamento é simulado; ambos continuam como limitação.
\item \textbf{Não coberto:} teste de carga e concorrência (por exemplo, dois clientes comprando a última unidade ao mesmo tempo), compatibilidade com outros navegadores (a interface foi testada em Chromium), acessibilidade, uso em celular e testes de invasão.
\item \textbf{Gateway de pagamento:} é simulado; o comportamento com um gateway real e com \textit{webhook} ainda não foi testado.
\item \textbf{Cancelamento:} cancelar um pedido não devolve as unidades ao estoque (limitação conhecida, documentada no manual de uso).
\item \textbf{Cobertura de código:} não foi medida com ferramenta; a cobertura é demonstrada pela rastreabilidade entre testes, histórias e regras.
\end{itemize}

\newpage
\appendix
\section{Casos de teste de back-end e banco de dados}

Legenda: \textbf{Refere-se a} indica a história (US), a regra de negócio (RN), o requisito não funcional (RNF) ou o tema (ACESSO, ESQUEMA, UTIL, ROBUSTEZ, VIEW) que o caso verifica.

\subsection*{JsonTest --- Conversor JSON (10 casos)}
\footnotesize
\begin{longtable}{|c|p{2.5cm}|p{8.4cm}|c|}
\hline
\rowcolor{gray!30} \textbf{N.º} & \textbf{Refere-se a} & \textbf{Caso de teste (comportamento esperado)} & \textbf{Result.} \\ \hline
\endhead
1 & UTIL & Escapa aspas, barra invertida e quebra de linha & \cellcolor{green!20}OK \\ \hline
2 & UTIL & Número decimal em campo inteiro gera IllegalArgumentException & \cellcolor{green!20}OK \\ \hline
3 & UTIL & Lê acentos e sequência \textbackslash\{\}u & \cellcolor{green!20}OK \\ \hline
4 & UTIL & Campo numérico obrigatório ausente gera IllegalArgumentException & \cellcolor{green!20}OK \\ \hline
5 & UTIL & Corpo vazio vira objeto vazio & \cellcolor{green!20}OK \\ \hline
6 & UTIL & Converte lista de ids separada por vírgula & \cellcolor{green!20}OK \\ \hline
7 & UTIL & Serializa mapa com texto, número e booleano & \cellcolor{green!20}OK \\ \hline
8 & UTIL & JSON malformado gera IllegalArgumentException & \cellcolor{green!20}OK \\ \hline
9 & UTIL & Serializa nulo e lista de mapas & \cellcolor{green!20}OK \\ \hline
10 & UTIL & Lê objeto simples com tipos misturados & \cellcolor{green!20}OK \\ \hline
\end{longtable}\normalsize
\subsection*{SenhaESessaoTest --- Hash de senha e sessões (5 casos)}
\footnotesize
\begin{longtable}{|c|p{2.5cm}|p{8.4cm}|c|}
\hline
\rowcolor{gray!30} \textbf{N.º} & \textbf{Refere-se a} & \textbf{Caso de teste (comportamento esperado)} & \textbf{Result.} \\ \hline
\endhead
11 & UTIL & Sessão: tokens diferentes para logins diferentes & \cellcolor{green!20}OK \\ \hline
12 & UTIL & Sessão: cria, busca e encerra & \cellcolor{green!20}OK \\ \hline
13 & UTIL & Hash é determinístico, tem 64 caracteres e muda com a senha & \cellcolor{green!20}OK \\ \hline
14 & UTIL & Hash SHA-256 de '123456' confere com o valor conhecido & \cellcolor{green!20}OK \\ \hline
15 & UTIL & Sessão: token nulo ou desconhecido não autentica & \cellcolor{green!20}OK \\ \hline
\end{longtable}\normalsize
\subsection*{UsuarioEEnderecoDaoTest --- Usuários, login e endereços (13 casos)}
\footnotesize
\begin{longtable}{|c|p{2.5cm}|p{8.4cm}|c|}
\hline
\rowcolor{gray!30} \textbf{N.º} & \textbf{Refere-se a} & \textbf{Caso de teste (comportamento esperado)} & \textbf{Result.} \\ \hline
\endhead
16 & ACESSO & E-mail repetido é rejeitado e não deixa registro órfão & \cellcolor{green!20}OK \\ \hline
17 & US07 & Lista os endereços do cliente & \cellcolor{green!20}OK \\ \hline
18 & US07 & Endereço inativo não é listado & \cellcolor{green!20}OK \\ \hline
19 & ACESSO & Senha errada ou e-mail desconhecido não autentica & \cellcolor{green!20}OK \\ \hline
20 & US07 & Salva novo endereço válido & \cellcolor{green!20}OK \\ \hline
21 & ACESSO & E-mail com formato inválido é rejeitado pelo banco & \cellcolor{green!20}OK \\ \hline
22 & ACESSO & Senha fica gravada apenas como hash & \cellcolor{green!20}OK \\ \hline
23 & ACESSO & Login de administrador devolve o perfil ADMINISTRADOR & \cellcolor{green!20}OK \\ \hline
24 & ACESSO & Usuário inativo não consegue entrar & \cellcolor{green!20}OK \\ \hline
25 & US07, RN07 & CEP inválido e UF em minúsculas são rejeitados pelo banco & \cellcolor{green!20}OK \\ \hline
26 & ACESSO & Login de cliente com senha correta & \cellcolor{green!20}OK \\ \hline
27 & US07, RN07 & Campos obrigatórios do endereço não podem faltar & \cellcolor{green!20}OK \\ \hline
28 & ACESSO & Cadastro cria usuário e cliente e permite login & \cellcolor{green!20}OK \\ \hline
\end{longtable}\normalsize
\subsection*{CupcakeDaoTest --- Vitrine, filtro e produtos (15 casos)}
\footnotesize
\begin{longtable}{|c|p{2.5cm}|p{8.4cm}|c|}
\hline
\rowcolor{gray!30} \textbf{N.º} & \textbf{Refere-se a} & \textbf{Caso de teste (comportamento esperado)} & \textbf{Result.} \\ \hline
\endhead
29 & US13 & Produto inativado sai da vitrine mas continua na lista do dono & \cellcolor{green!20}OK \\ \hline
30 & US12 & Cadastrar produto grava dados e categorias & \cellcolor{green!20}OK \\ \hline
31 & US03 & Detalhes trazem descrição, ingredientes, preço e categorias & \cellcolor{green!20}OK \\ \hline
32 & US12, RN12 & Nome de produto repetido é rejeitado & \cellcolor{green!20}OK \\ \hline
33 & US13, RN13 & Editar produto inexistente gera erro & \cellcolor{green!20}OK \\ \hline
34 & US02 & Filtro não mostra produto inativo da categoria & \cellcolor{green!20}OK \\ \hline
35 & US02 & Lista de categorias vem ordenada por nome & \cellcolor{green!20}OK \\ \hline
36 & US01 & Vitrine mostra os produtos em destaque primeiro & \cellcolor{green!20}OK \\ \hline
37 & US12 & Falha ao gravar categoria desfaz o cadastro inteiro (transação) & \cellcolor{green!20}OK \\ \hline
38 & US03 & Produto inexistente retorna nulo & \cellcolor{green!20}OK \\ \hline
39 & US02 & Filtro por categoria retorna apenas produtos da categoria & \cellcolor{green!20}OK \\ \hline
40 & US02, RN02 & Categoria inexistente retorna lista vazia & \cellcolor{green!20}OK \\ \hline
41 & US13, RN03 & Editar produto altera preço, estoque e categorias (dado atualizado) & \cellcolor{green!20}OK \\ \hline
42 & US12, RN12 & Preço zero e estoque negativo são rejeitados pelo banco & \cellcolor{green!20}OK \\ \hline
43 & US01, RN01 & Vitrine lista somente produtos ativos & \cellcolor{green!20}OK \\ \hline
\end{longtable}\normalsize
\subsection*{CarrinhoDaoTest --- Carrinho de compras (11 casos)}
\footnotesize
\begin{longtable}{|c|p{2.5cm}|p{8.4cm}|c|}
\hline
\rowcolor{gray!30} \textbf{N.º} & \textbf{Refere-se a} & \textbf{Caso de teste (comportamento esperado)} & \textbf{Result.} \\ \hline
\endhead
44 & US05 & Remover item que não está no carrinho não gera erro & \cellcolor{green!20}OK \\ \hline
45 & US04 & Produto inexistente não pode ser adicionado & \cellcolor{green!20}OK \\ \hline
46 & US04 & Adicionar item cria o carrinho aberto e grava a quantidade & \cellcolor{green!20}OK \\ \hline
47 & US04 & Adicionar o mesmo produto de novo soma as quantidades & \cellcolor{green!20}OK \\ \hline
48 & US06 & Carrinho de um cliente não aparece para outro & \cellcolor{green!20}OK \\ \hline
49 & US06, RN06 & Carrinho calcula subtotal por item & \cellcolor{green!20}OK \\ \hline
50 & US06 & Carrinho vazio retorna lista vazia & \cellcolor{green!20}OK \\ \hline
51 & US04, RN05 & Quantidade negativa é rejeitada & \cellcolor{green!20}OK \\ \hline
52 & US04, RN04 & Quantidade zero é rejeitada e nada é gravado & \cellcolor{green!20}OK \\ \hline
53 & US05 & Remover item tira o produto e o total é recalculado & \cellcolor{green!20}OK \\ \hline
54 & US04 & Produto inativo não pode ser adicionado & \cellcolor{green!20}OK \\ \hline
\end{longtable}\normalsize
\subsection*{PedidoDaoTest --- Pedido, pagamento, status e dashboard (29 casos)}
\footnotesize
\begin{longtable}{|c|p{2.5cm}|p{8.4cm}|c|}
\hline
\rowcolor{gray!30} \textbf{N.º} & \textbf{Refere-se a} & \textbf{Caso de teste (comportamento esperado)} & \textbf{Result.} \\ \hline
\endhead
55 & US14, RN11 & Status inexistente e pedido inexistente geram erro & \cellcolor{green!20}OK \\ \hline
56 & US10 & Cliente não consegue pagar o pedido de outro cliente & \cellcolor{green!20}OK \\ \hline
57 & US10 & Pagamento recusado mantém o pedido aguardando e permite nova tentativa & \cellcolor{green!20}OK \\ \hline
58 & US11 & Pedido nasce com o status inicial registrado no histórico & \cellcolor{green!20}OK \\ \hline
59 & US09 & Endereço de entrega é copiado para o pedido & \cellcolor{green!20}OK \\ \hline
60 & US09, RN09 & Estoque insuficiente impede o pedido e nada é alterado & \cellcolor{green!20}OK \\ \hline
61 & US09, RN09 & Sem carrinho aberto não é possível confirmar & \cellcolor{green!20}OK \\ \hline
62 & US10, RN10 & Pagamento aprovado muda o pedido para PAGO & \cellcolor{green!20}OK \\ \hline
63 & US15, RN15 & Vendas só contam pedidos já pagos e refletem o dado atual & \cellcolor{green!20}OK \\ \hline
64 & US08, RN08 & Forma de pagamento inválida é rejeitada e nada é gravado & \cellcolor{green!20}OK \\ \hline
65 & US09 & Confirmar cria pedido aguardando pagamento com total = itens + taxa & \cellcolor{green!20}OK \\ \hline
66 & US09, RN07 & Endereço de outro cliente ou inexistente é recusado & \cellcolor{green!20}OK \\ \hline
67 & US11 & Cliente vê apenas os seus pedidos, com status e pagamento & \cellcolor{green!20}OK \\ \hline
68 & US09 & Depois de confirmar, o cliente pode montar um novo carrinho & \cellcolor{green!20}OK \\ \hline
69 & US14, RN14 & Pedido entregue não pode mais mudar de status & \cellcolor{green!20}OK \\ \hline
70 & US10 & Não é possível pagar um pedido que já foi pago & \cellcolor{green!20}OK \\ \hline
71 & US09, RN09 & Carrinho vazio não gera pedido & \cellcolor{green!20}OK \\ \hline
72 & US14 & Histórico guarda quem alterou o status & \cellcolor{green!20}OK \\ \hline
73 & US09 & Preço do item fica gravado: mudar o preço depois não altera o pedido & \cellcolor{green!20}OK \\ \hline
74 & US09, RN09 & Produto inativado depois de entrar no carrinho impede o pedido & \cellcolor{green!20}OK \\ \hline
75 & US11 & Pedidos aparecem do mais novo para o mais antigo & \cellcolor{green!20}OK \\ \hline
76 & US09 & Confirmar fecha o carrinho e esvazia a lista & \cellcolor{green!20}OK \\ \hline
77 & US14, RN14 & Pedido pode ser cancelado enquanto aguarda pagamento & \cellcolor{green!20}OK \\ \hline
78 & US09 & Confirmar baixa o estoque dos produtos & \cellcolor{green!20}OK \\ \hline
79 & US15 & Pedido cancelado não entra nas vendas & \cellcolor{green!20}OK \\ \hline
80 & US14, RN14 & Pular etapa do fluxo é rejeitado & \cellcolor{green!20}OK \\ \hline
81 & US15 & Mais vendidos vem ordenado por unidades & \cellcolor{green!20}OK \\ \hline
82 & US14, RN14 & Fluxo completo de status: PAGO, EM\_PREPARO, ENVIADO e ENTREGUE & \cellcolor{green!20}OK \\ \hline
83 & US15 & Relatório de estoque classifica BAIXO, OK e ESGOTADO & \cellcolor{green!20}OK \\ \hline
\end{longtable}\normalsize
\subsection*{BancoDeDadosTest --- Restrições, triggers, views e permissões (13 casos)}
\footnotesize
\begin{longtable}{|c|p{2.5cm}|p{8.4cm}|c|}
\hline
\rowcolor{gray!30} \textbf{N.º} & \textbf{Refere-se a} & \textbf{Caso de teste (comportamento esperado)} & \textbf{Result.} \\ \hline
\endhead
84 & US04, ESQUEMA & Um cliente só pode ter um carrinho aberto & \cellcolor{green!20}OK \\ \hline
85 & RNF08 & Papel da aplicação não pode apagar pedidos nem pagamentos & \cellcolor{green!20}OK \\ \hline
86 & ACESSO & Senhas de demonstração estão guardadas como hash & \cellcolor{green!20}OK \\ \hline
87 & US01, US15 & Views de vitrine e estoque respondem com os dados iniciais & \cellcolor{green!20}OK \\ \hline
88 & ESQUEMA & Subtotal do item do pedido é calculado pelo banco & \cellcolor{green!20}OK \\ \hline
89 & ESQUEMA & Valores monetários negativos são rejeitados & \cellcolor{green!20}OK \\ \hline
90 & US10, RN10 & Pagamento processado precisa ter data de processamento & \cellcolor{green!20}OK \\ \hline
91 & ESQUEMA & Nome de categoria é único & \cellcolor{green!20}OK \\ \hline
92 & ESQUEMA & Produto com pedidos não pode ser apagado & \cellcolor{green!20}OK \\ \hline
93 & RN14 & Tabela de transições tem as 7 regras do fluxo & \cellcolor{green!20}OK \\ \hline
94 & US10, ESQUEMA & Um pedido não pode ter dois pagamentos aprovados & \cellcolor{green!20}OK \\ \hline
95 & US07, ESQUEMA & Só um endereço principal ativo por cliente & \cellcolor{green!20}OK \\ \hline
96 & ESQUEMA & Administrador não pode ser cadastrado como cliente & \cellcolor{green!20}OK \\ \hline
\end{longtable}\normalsize
\subsection*{ApiTest --- Rotas, perfis, validações e fluxos completos (35 casos)}
\footnotesize
\begin{longtable}{|c|p{2.5cm}|p{8.4cm}|c|}
\hline
\rowcolor{gray!30} \textbf{N.º} & \textbf{Refere-se a} & \textbf{Caso de teste (comportamento esperado)} & \textbf{Result.} \\ \hline
\endhead
97 & ROBUSTEZ & Método HTTP não permitido devolve 405; rota inexistente devolve 404 & \cellcolor{green!20}OK \\ \hline
98 & US03 & Detalhes do produto: 200 se existe e ativo, 404 se não & \cellcolor{green!20}OK \\ \hline
99 & ACESSO & Rotas protegidas exigem login (401) e rejeitam token inventado & \cellcolor{green!20}OK \\ \hline
100 & ACESSO & Login válido devolve token, nome e perfil & \cellcolor{green!20}OK \\ \hline
101 & US12 & Dono cadastra produto novo e ele aparece na vitrine & \cellcolor{green!20}OK \\ \hline
102 & ACESSO & Cadastro rejeita e-mail repetido, senha curta e telefone inválido & \cellcolor{green!20}OK \\ \hline
103 & US13, RN03 & Dono edita preço e estoque (detalhes mostram o dado novo); desativar tira da vitrine & \cellcolor{green!20}OK \\ \hline
104 & US04, RN04 & Quantidade fora de 1 a 99 é recusada com 400 & \cellcolor{green!20}OK \\ \hline
105 & US14, RN11 & Status inexistente, ausente ou pedido inexistente devolvem 400 & \cellcolor{green!20}OK \\ \hline
106 & ACESSO & Cadastro de cliente novo e login em seguida & \cellcolor{green!20}OK \\ \hline
107 & US09, RN09 & Confirmar com carrinho vazio ou sem endereço devolve 400 & \cellcolor{green!20}OK \\ \hline
108 & US02, RN02 & Categoria inválida gera erro 400 & \cellcolor{green!20}OK \\ \hline
109 & ROBUSTEZ & JSON malformado devolve 400 (e não erro 500) & \cellcolor{green!20}OK \\ \hline
110 & US01 & Vitrine é pública e lista os 5 cupcakes ativos & \cellcolor{green!20}OK \\ \hline
111 & US11 & Cliente só vê os próprios pedidos & \cellcolor{green!20}OK \\ \hline
112 & US10, RN10 & Pagamento recusado e depois aprovado leva o pedido a PAGO & \cellcolor{green!20}OK \\ \hline
113 & US01, RN01 & Produto inativo não aparece na vitrine da API & \cellcolor{green!20}OK \\ \hline
114 & ACESSO & Logout invalida o token & \cellcolor{green!20}OK \\ \hline
115 & ACESSO & Cliente não acessa o painel do dono (403) e o dono não usa o carrinho (403) & \cellcolor{green!20}OK \\ \hline
116 & US07 & Cadastrar endereço válido e listá-lo & \cellcolor{green!20}OK \\ \hline
117 & VIEW & Páginas estáticas são servidas; fora da pasta public dá 404 & \cellcolor{green!20}OK \\ \hline
118 & ACESSO & Login inválido devolve 401; campos ausentes devolvem 400 & \cellcolor{green!20}OK \\ \hline
119 & US12, RN12 & Cadastro recusa campos obrigatórios vazios, preço zero e nome repetido & \cellcolor{green!20}OK \\ \hline
120 & US14, RN14 & Dono muda o status seguindo o fluxo e é barrado ao pular etapas & \cellcolor{green!20}OK \\ \hline
121 & RNF01, RNF02 & Vitrine e filtro respondem em menos de 1 segundo & \cellcolor{green!20}OK \\ \hline
122 & US13, RN13 & Editar produto inexistente devolve 400 & \cellcolor{green!20}OK \\ \hline
123 & ROBUSTEZ & Texto com aspas e acentos no nome do produto volta íntegro no JSON & \cellcolor{green!20}OK \\ \hline
124 & US09 & Confirmar pedido: 201, estoque baixa e carrinho esvazia & \cellcolor{green!20}OK \\ \hline
125 & US08, RN08 & Forma de pagamento inválida devolve 400 & \cellcolor{green!20}OK \\ \hline
126 & US10 & Pedido já pago não pode ser pago de novo; pedido alheio também não & \cellcolor{green!20}OK \\ \hline
127 & US07, RN07 & Endereço com CEP, UF ou campos obrigatórios inválidos é recusado & \cellcolor{green!20}OK \\ \hline
128 & US15 & Dashboard devolve vendas, mais vendidos e estoque & \cellcolor{green!20}OK \\ \hline
129 & US02 & Filtro por categoria retorna só os produtos dela & \cellcolor{green!20}OK \\ \hline
130 & US04, US06 & Adicionar ao carrinho devolve itens e total atualizados & \cellcolor{green!20}OK \\ \hline
131 & US05 & Remover item do carrinho atualiza o total & \cellcolor{green!20}OK \\ \hline
\end{longtable}\normalsize

\newpage
\section{Casos de teste de interface}

\footnotesize
\begin{longtable}{|c|c|p{9.4cm}|c|}
\hline
\rowcolor{gray!30} \textbf{Cód.} & \textbf{Refere-se a} & \textbf{Caso de teste de interface (comportamento esperado)} & \textbf{Result.} \\ \hline
\endhead
UI01 & ACESSO & Ao abrir, aparece somente a tela de login (loja e vitrine ocultas) & \cellcolor{green!20}OK \\ \hline
UI02 & ACESSO & Login com senha errada mostra mensagem de erro e continua no login & \cellcolor{green!20}OK \\ \hline
UI03 & ACESSO & Cadastro de conta nova volta ao login com o e-mail preenchido & \cellcolor{green!20}OK \\ \hline
UI04 & ACESSO & Cadastro com senha curta mostra erro & \cellcolor{green!20}OK \\ \hline
UI05 & US01 & Após o login aparece a vitrine com 5 cupcakes e some a tela de login & \cellcolor{green!20}OK \\ \hline
UI06 & US01 & Produtos em destaque aparecem primeiro & \cellcolor{green!20}OK \\ \hline
UI07 & US02 & Filtro 'Vegano' mostra 1 cupcake e 'Todas' volta aos 5 & \cellcolor{green!20}OK \\ \hline
UI08 & US03 & Detalhes abre janela com ingredientes e preço & \cellcolor{green!20}OK \\ \hline
UI09 & US04 & Adicionar mostra aviso e atualiza o contador do carrinho & \cellcolor{green!20}OK \\ \hline
UI10 & US06 & Aba lateral abre o carrinho com 2 itens e total R\$ 31,50 & \cellcolor{green!20}OK \\ \hline
UI11 & US05 & Remover item atualiza lista, total e contador & \cellcolor{green!20}OK \\ \hline
UI12 & US05 & Carrinho vazio mostra mensagem e desabilita Finalizar pedido & \cellcolor{green!20}OK \\ \hline
UI13 & ACESSO & Clicar fora da aba lateral (ou no X) fecha o painel & \cellcolor{green!20}OK \\ \hline
UI14 & US09 & Finalizar pedido: escolher endereço da Maria e confirmar abre o pagamento & \cellcolor{green!20}OK \\ \hline
UI15 & US10 & Pagamento recusado avisa o cliente; nova tentativa é aprovada & \cellcolor{green!20}OK \\ \hline
UI16 & US11 & Meus pedidos mostra o pedido com status PAGO e pagamento APROVADO & \cellcolor{green!20}OK \\ \hline
UI17 & ACESSO & Recarregar a página mantém o usuário logado & \cellcolor{green!20}OK \\ \hline
UI18 & ACESSO & Cliente logado não vê o painel do dono em admin.html & \cellcolor{green!20}OK \\ \hline
UI19 & ACESSO & Sair devolve à tela de login & \cellcolor{green!20}OK \\ \hline
UI20 & US07 & Sem endereço, a janela de finalização já abre o formulário de novo endereço & \cellcolor{green!20}OK \\ \hline
UI21 & US07 & CEP inválido mostra erro e o endereço não é salvo & \cellcolor{green!20}OK \\ \hline
UI22 & US07 & Endereço válido é salvo e já fica selecionado & \cellcolor{green!20}OK \\ \hline
UI23 & US09 & Pedido confirmado fica 'aguardando pagamento' com botão Pagar agora & \cellcolor{green!20}OK \\ \hline
UI24 & ACESSO & Administrador que entra pela tela de login vai direto ao painel do dono & \cellcolor{green!20}OK \\ \hline
UI25 & US15 & Dashboard mostra vendas, mais vendidos e estoque & \cellcolor{green!20}OK \\ \hline
UI26 & US14 & Dono muda o status de PAGO para EM\_PREPARO & \cellcolor{green!20}OK \\ \hline
UI27 & US14 & Pular etapa (EM\_PREPARO para ENTREGUE) mostra a mensagem de erro & \cellcolor{green!20}OK \\ \hline
UI28 & US12 & Dono cadastra produto novo e ele aparece na lista de produtos & \cellcolor{green!20}OK \\ \hline
UI29 & US12 & Cadastro com preço vazio é barrado na tela & \cellcolor{green!20}OK \\ \hline
UI30 & US13 & Dono edita o estoque de um produto (repor estoque) & \cellcolor{green!20}OK \\ \hline
UI31 & SEGURANÇA & Nome de produto com HTML aparece como texto (sem executar código) & \cellcolor{green!20}OK \\ \hline
UI32 & US04, RN09 & Produto esgotado aparece com a marca "Esgotado" e o botão de compra desabilitado (D4) & \cellcolor{green!20}OK \\ \hline
UI33 & US04 & Escolher a quantidade na vitrine adiciona várias unidades de uma vez (D5) & \cellcolor{green!20}OK \\ \hline
UI34 & US05 & No carrinho, os botões + e − e o campo numérico alteram a quantidade e o total (D5) & \cellcolor{green!20}OK \\ \hline
UI35 & RN09 & Quantidade acima do estoque no carrinho é recusada com mensagem clara e o valor volta ao anterior (D4) & \cellcolor{green!20}OK \\ \hline
UI36 & US04, RN09 & Adicionar mais do que o estoque pela vitrine mostra o erro e mantém o carrinho (D4) & \cellcolor{green!20}OK \\ \hline
UI37 & US15 & Dashboard do dono atualiza sozinho, sem recarregar a página (D6) & \cellcolor{green!20}OK \\ \hline
UI38 & US15 & Botão "Atualizar agora" recarrega o dashboard sob demanda & \cellcolor{green!20}OK \\ \hline
\end{longtable}\normalsize

\end{document}
